\begin{figure}[t]
\centering
\begin{minipage}[t]{0.58\linewidth}
\centering
\begin{tikzpicture}
\begin{axis}[
  width=0.95\linewidth,
  height=0.67\linewidth,
  xmin=0,xmax=3,
  ymin=0,ymax=0.41,
  xlabel={$\lambda m$},
  ylabel={$r m$},
  xtick={0,1,2,3},
  ytick={0,0.1,0.2,0.3,0.367879},
  yticklabels={$0$,$0.1$,$0.2$,$0.3$,$1/e$},
  tick align=outside,
]
\addplot[black,thick,domain=0:3,samples=120] {x*exp(-x)};
\addplot[only marks,mark=*,mark size=1.8pt,black] coordinates {(1,0.3678794412)};
\draw[densely dashed] (axis cs:1,0) -- (axis cs:1,0.3678794412);
\node[anchor=south east,font=\scriptsize,align=right] at (axis cs:2.88,0.305) {identical for all recovery laws\\with mean $m$};
\end{axis}
\end{tikzpicture}
\end{minipage}\hfill
\begin{minipage}[t]{0.38\linewidth}
\centering
\begin{tikzpicture}
\begin{axis}[
  width=0.95\linewidth,
  height=1.02\linewidth,
  ybar,
  bar width=13pt,
  symbolic x coords={deterministic,exponential},
  xtick=data,
  x tick label style={rotate=25,anchor=east,font=\scriptsize},
  ymin=0,ymax=0.08,
  ylabel={$\mathcal G_{\rm DC}$ at $\lambda m=1$},
  ytick={0,0.02,0.04,0.06,0.08},
  nodes near coords,
  nodes near coords style={font=\scriptsize,/pgf/number format/fixed,/pgf/number format/precision=4},
  tick align=outside,
]
\addplot[fill=white,draw=black,thick] coordinates {(deterministic,0) (exponential,0.06915579)};
\end{axis}
\end{tikzpicture}
\end{minipage}
\caption{The conventional saturation curve does not determine static information retention. Left: every iid Type-II recovery law with mean $m$ has the same registered-rate curve $r=\lambda e^{-\lambda m}$ and the same maximum at $\lambda m=1$. Right: at that common operating point deterministic recovery has exactly zero stationary Fisher retention, whereas exponential recovery retains $\mathcal G_{\rm DC}\simeq0.06916$.}
\label{fig:sameCurve}
\end{figure}
